\documentclass{situationist-preprint}

\title{The Reality of Rights}
\msid{POL-18-a}
\msfamily{POL-18}
\msdiscipline{Political Philosophy}
\mstype{Principal manuscript}
\mskeywords{rights; natural rights; Bentham; claim-rights; Hohfeld; exclusionary reasons; directed duties; legal positivism}

\begin{document}
\maketitle

\begin{abstract}
Bentham called natural rights simple nonsense, and natural and imprescriptible rights rhetorical nonsense, nonsense upon stilts. His argument was that a right is the creature of a law, a law is a command backed by a sanction, and where there is no lawgiver there is no right, only a wish that there were one. We argue that rights exist and are real normative constraints. A claim-right of one person against another is a structure of reasons: a reason for the second not to act in a certain way, which is owed to the first, which the first may waive, and which excludes from the second's deliberation certain reasons for acting that would otherwise count. Given that reasons are real and primitive, rights so understood are real. We show that some such structures obtain independently of any law, and that they cannot be reduced to facts about the utility of compliance, since waiver changes what may be done without changing any outcome. Bentham's argument rests on the command theory of normativity, which Hart showed fails even for legal rights. His practical worry, that natural rights are anarchical, is answered by their defeasibility.
\end{abstract}

\section{The problem}

The Declaration of the Rights of Man and of the Citizen of 1789 proclaimed that men are born and remain free and equal in rights, and that the aim of every political association is the preservation of the natural and imprescriptible rights of man. \citet{Bentham1843} replied, in a text written in the 1790s and published only later, that natural rights are simple nonsense, and natural and imprescriptible rights rhetorical nonsense, nonsense upon stilts. The phrase names a position that remains live: that there are no rights except those conferred by positive law or convention, and that talk of rights beyond these is confused or disguises a wish, a demand or a claim about utility. \citet[ch.~6]{MacIntyre1981} put it more starkly: belief in natural or human rights is of a piece with belief in witches and unicorns, since every attempt to give good reasons for believing in them has failed.

Bentham's argument has two parts \citep{Waldron1987}. The first is conceptual: a right is correlative to a duty; a duty is what one is commanded to do on pain of a sanction; a command requires a commander; so a right requires a lawgiver, and natural rights, which are supposed to exist without one, are a contradiction. The second is practical: declarations of natural and imprescriptible rights are anarchical.

Our thesis is that rights exist and are real normative constraints. The account of rights is the one MTE-11-a sketches in its application to rights, and the argument depends on its thesis that reasons are primitive and real.

\section{Framework}

Of the four Hohfeldian elements, claims, privileges, powers and immunities \citep{Hohfeld1919}, we are concerned with claims, which are correlative to duties and are the target of Bentham's argument.

\begin{definition}[Claim-right]\label{d:claim}
A person $x$ has a \emph{claim} against a person $y$ that $y$ not $\phi$ if and only if there is a reason for $y$ not to $\phi$ such that
\begin{enumerate}[label=(\alph*)]
\item it is \emph{owed to} $x$: if $y$ $\phi$s without justification, $y$ wrongs $x$, and not merely acts wrongly;
\item it is \emph{waivable by} $x$: if $x$ validly consents to $y$'s $\phi$-ing, the reason lapses;
\item it is \emph{exclusionary}: it excludes from $y$'s deliberation certain reasons for $\phi$-ing that would otherwise count, such as the benefit to others of a modest infringement.
\end{enumerate}
\end{definition}

Definition~\ref{d:claim} combines three features that rival theories of rights have each emphasized. Clause (a) is the directedness that \citet[ch.~1]{Thomson1990} and \citet{Darwall2006} place at the centre. Clause (b) is the control that the will theory, defended by \citet[ch.~7]{Hart1982}, takes to be essential. Clause (c) is the structure that \citet[ch.~3]{Nozick1974} called a side constraint and \citet{Dworkin1977} a trump, and that \citet{Raz1999} analyses as an exclusionary reason. Not every right has all three features (those of infants are not waivable by their holders), but the central cases do, and a structure of reasons with all three is a right.

\begin{definition}[Normative constraint]\label{d:constraint}
A \emph{normative constraint} on an agent is a consideration that limits the options the agent may permissibly take, independently of the agent's goals and without being outweighed by any aggregate of benefits below a threshold.
\end{definition}

\begin{definition}[Reality]\label{d:real}
A normative item is \emph{real} if its obtaining is a normative fact that is not constituted by any fact about what a legal system or community commands, recognizes or enforces.
\end{definition}

\begin{assumption}[Reasons first]\label{a:reasons}
Reasons are real and primitive, and the other normative notions are defined in terms of them.
\end{assumption}

Assumption~\ref{a:reasons} is the conclusion of MTE-11. It says that facts about reasons are not constituted by facts about commands or conventions, which is what Definition~\ref{d:real} requires of real normative items, and that other normative notions, including rights, are relations among reasons.

\section{The argument}

\begin{lemma}\label{l:structure}
A claim-right is a normative constraint.
\end{lemma}

\begin{proof}
Let $x$ have a claim against $y$ that $y$ not $\phi$. By clause (c) of Definition~\ref{d:claim}, the reason not to $\phi$ excludes from $y$'s deliberation reasons for $\phi$-ing such as modest benefits to others, so it is not outweighed by any aggregate of such benefits; and it limits $y$'s permissible options whatever $y$'s goals are, since a reason owed to another does not depend on the goals of the person who owes it. By Definition~\ref{d:constraint}, it is a normative constraint.
\end{proof}

\begin{lemma}[Independence of law]\label{l:indep}
Some claim-rights obtain in the absence of any law that confers them.
\end{lemma}

\begin{proof}
Consider a regime whose laws permit officials to torture suspects to extract information, and an official who tortures a suspect under that law. Two judgements about the case are as secure as any in ethics: that the official acts wrongly, and that he wrongs the suspect, to whom apology and redress are owed. The second is the judgement that he has violated a reason owed to the suspect, which is clause (a) of Definition~\ref{d:claim}. The reason is of a kind consent can waive, since the same infliction of pain, validly consented to as part of a medical procedure, would wrong no one; this is clause (b). And the information extracted, even if it benefits many, does not make the torture permissible; this is clause (c). So the suspect has a claim against the official, and it obtains though the law confers no such right and in fact denies it.
\end{proof}

\begin{lemma}[Irreducibility to utility]\label{l:irred}
Whether a claim-right obtains is not settled by facts about the utility of compliance and non-compliance.
\end{lemma}

\begin{proof}
Compare two cases that agree in every outcome. In the first, a surgeon performs an operation on a patient who has consented; in the second, the same surgeon performs the same operation, with the same results, on a patient who has not. The outcomes, including the patient's health, pain and satisfaction afterwards, may be identical; a patient may be glad of an operation she did not authorize. Yet the first operation is permissible and the second wrongs the patient. What distinguishes the cases is the waiver of clause (b) of Definition~\ref{d:claim}, which changes the reasons that bear on the surgeon without changing any outcome. A fact that is not settled by the outcomes is not settled by their utility.
\end{proof}

\begin{theorem}\label{t:main}
Rights exist, and they are real normative constraints.
\end{theorem}

\begin{proof}
By Lemma~\ref{l:indep}, some claim-rights obtain. By Definition~\ref{d:claim}, a claim-right is a structure of reasons, and by Assumption~\ref{a:reasons} facts about reasons are real normative facts, not constituted by any law or convention. Lemma~\ref{l:indep} confirms this in the case of the tortured suspect, whose right is not constituted by the law, which denies it. By Lemma~\ref{l:structure}, claim-rights are normative constraints. Lemma~\ref{l:irred} shows that the constraint is not a disguised claim about utility, so that Bentham's suggested paraphrase of rights talk does not eliminate it.
\end{proof}

\begin{remark}\label{r:eth02}
ETH-02 argues that every deontological theory is extensionally equivalent to a consequentialist theory with a suitable value function. This does not tell against Theorem~\ref{t:main}: the value function that reproduces a theory with rights must be agent-relative and sensitive to waiver, and so encodes the structure of Definition~\ref{d:claim} in its content. The equivalence shows that the constraints can be represented in either notation, not that there are none to represent.
\end{remark}

\section{Where Bentham's argument fails}

Bentham's conceptual argument assumes that a duty is what one is commanded to do on pain of a sanction. That is the command theory of normativity, and \citet[chs.~2--4]{Hart1961} showed that it fails even for law. A gunman who orders a clerk to hand over money obliges him to comply but imposes no obligation on him. Legal obligations are standards of conduct accepted from an internal point of view by those who operate the system, and the sanction responds to their breach rather than constituting them. If even legal duties are not commands backed by sanctions, the inference from the absence of a commander to the absence of a duty has no support, and Bentham's own account of legal rights fails with it.

Bentham also insisted that a reason for wishing there were a right is not a right, as hunger is not bread. The point is correct and harmless: a right in the sense of Definition~\ref{d:claim} is not a reason for wishing that a law conferred one, but a reason, owed to the right-holder now, that bears on another's action. The official has a reason not to torture the suspect, owed to him, which the law ignores.

Bentham's practical argument is that declarations of natural and imprescriptible rights are anarchical, since they invite every person to resist any law he judges to infringe them. The argument has force against the drafting of 1789, which stated rights without limits, but not against the existence of rights. Rights in the sense of Definition~\ref{d:claim} have thresholds: clause (c) excludes reasons below a threshold, not all reasons, and \citet[ch.~3]{Thomson1990} has shown how a right may be permissibly \emph{infringed} when the stakes are high enough without being \emph{violated}. And who judges whether a right has been infringed is a question about institutions: POL-02 argues that a legitimate state is one that can be justified to every reasonable person subject to it, and a state that protects rights through settled law and independent courts gives each person a judge other than herself, which is the remedy for the anarchy Bentham feared.

\section{Objections}

\begin{objection}[Queerness]
Rights so understood are odd entities: reasons that are owed to persons, that persons can cancel by fiat, and that silence other reasons. Nothing in the natural world is like that. Bentham's suspicion of fictitious entities was a naturalist suspicion, and it applies.
\end{objection}

\begin{reply}
The objection is an instance of Mackie's argument from queerness, which MTE-05 argues proves too much, since it would make mathematics false as well. That a reason can be owed to a person, and that her consent can change it, are relations among reasons, and if reasons exist, as Assumption~\ref{a:reasons} holds, there is no further difficulty in their having these relations.
\end{reply}

\begin{objection}[Witches and unicorns]
\citet[ch.~6]{MacIntyre1981} argues that every attempt to give good reasons for believing in rights has failed, and that the concept arose late, in a particular culture. A concept with so local a history is unlikely to answer to anything real.
\end{objection}

\begin{reply}
The history of a concept does not settle whether it applies; the late articulation of the concept of a gene does not show that heredity began in the nineteenth century. As for reasons, Lemma~\ref{l:indep} gives one: the judgement that the tortured suspect is wronged, and not merely affected by a wrong, does not depend on any theory of rights. \citet{Feinberg1970} made the point by imagining Nowheresville, a world with duties and benevolence but no rights, whose inhabitants lack the standing to claim what is owed to them; the lack is felt as a loss.
\end{reply}

\begin{objection}[Rights as rules of thumb]
A utilitarian can accept that we should think as if there were rights, since respecting them produces better outcomes than case-by-case calculation. Rights are then useful fictions, and Bentham was right about their ultimate status.
\end{objection}

\begin{reply}
On a two-level theory such as \citet{Hare1981}'s, the critical level can in principle override the intuitive one, and rights are respected only because we cannot reliably identify the cases where violation is best. Lemma~\ref{l:irred} tells against this: the unauthorized operation wrongs the patient even when we know its outcome is as good as the authorized one. A theory whose ultimate facts concern outcomes can represent the waiver only by building it into the value of outcomes, and then, as Remark~\ref{r:eth02} notes, the right is encoded, not eliminated.
\end{reply}

\begin{objection}[Positivism about rights]
Rights may be real without being natural: they may be creatures of social convention, real as money is real. MET-20 argues that social kinds are no less real than natural kinds. Then rights are real, as the thesis says, and Bentham was right that they are made by human practice.
\end{objection}

\begin{reply}
We agree that many rights are conventional: property rights in their detail, contractual rights, the rights of office. Their reality is that of social kinds, which MET-20 defends. But Lemma~\ref{l:indep} shows that some rights obtain where the conventions deny them: the tortured suspect's right is not a social fact of his society, whose social facts are against it. If all rights were conventional, there would be no standpoint from which to say that a society permitting torture wrongs its victims, rather than that it has different conventions. Conventional rights are, in the central cases, the institutional specification of rights that exist independently.
\end{reply}

\section{What the result does not establish}

The argument does not establish a list of rights, the threshold at which a right may be infringed, or whether any right is absolute. It does not decide between the will and the interest theories, since Definition~\ref{d:claim} concerns the central cases. And it does not vindicate every declaration of rights; Bentham's criticism of particular declarations may be sound where his denial of rights is not.

\section{Conclusion}

Rights are not nonsense upon stilts. A claim-right is a reason owed to its holder, waivable by her, and exclusionary; if reasons are real, so are rights. Some rights obtain where no law confers them, and consent changes what may be done without changing any outcome. Bentham's denial rested on the command theory, which fails even for law, and his fear of anarchy is answered by protecting rights through institutions.

\begin{thebibliography}{99}\small
\bibitem[Bentham(1843)]{Bentham1843} Bentham, J. (1843). Anarchical fallacies. In J.~Bowring (ed.), \emph{The Works of Jeremy Bentham}, Vol.~2. William Tait. Written in the 1790s.
\bibitem[Darwall(2006)]{Darwall2006} Darwall, S. (2006). \emph{The Second-Person Standpoint: Morality, Respect, and Accountability}. Harvard University Press.
\bibitem[Dworkin(1977)]{Dworkin1977} Dworkin, R. (1977). \emph{Taking Rights Seriously}. Harvard University Press.
\bibitem[Feinberg(1970)]{Feinberg1970} Feinberg, J. (1970). The nature and value of rights. \emph{The Journal of Value Inquiry}, 4(4), 243--260.
\bibitem[Hare(1981)]{Hare1981} Hare, R.~M. (1981). \emph{Moral Thinking: Its Levels, Method, and Point}. Clarendon Press.
\bibitem[Hart(1961)]{Hart1961} Hart, H.~L.~A. (1961). \emph{The Concept of Law}. Clarendon Press.
\bibitem[Hart(1982)]{Hart1982} Hart, H.~L.~A. (1982). \emph{Essays on Bentham: Studies in Jurisprudence and Political Theory}. Clarendon Press.
\bibitem[Hohfeld(1919)]{Hohfeld1919} Hohfeld, W.~N. (1919). \emph{Fundamental Legal Conceptions as Applied in Judicial Reasoning}. Ed. W.~W. Cook. Yale University Press.
\bibitem[MacIntyre(1981)]{MacIntyre1981} MacIntyre, A. (1981). \emph{After Virtue: A Study in Moral Theory}. Duckworth.
\bibitem[Nozick(1974)]{Nozick1974} Nozick, R. (1974). \emph{Anarchy, State, and Utopia}. Basic Books.
\bibitem[Raz(1999)]{Raz1999} Raz, J. (1999). \emph{Practical Reason and Norms}. Oxford University Press. First published 1975.
\bibitem[Thomson(1990)]{Thomson1990} Thomson, J.~J. (1990). \emph{The Realm of Rights}. Harvard University Press.
\bibitem[Waldron(1987)]{Waldron1987} Waldron, J. (ed.) (1987). \emph{\enquote{Nonsense upon Stilts}: Bentham, Burke and Marx on the Rights of Man}. Methuen.
\end{thebibliography}

\end{document}
